% Space asked for inside a paragraph: \\[len], \vspace and \smallskip mid-paragraph, \vadjust pre, and a box and a rule in a \vadjust. TeX puts it below (or above) the line it ends up in.
\documentclass{article}
% one page with nothing but the text: checked against plain lualatex too (cases.ts)
\pagestyle{empty}
\begin{document}
Case A, plain text: the first line of the paragraph.\\[4pt]
The second line, after a newline with four points.\\
The third line, after a plain newline.\\[10pt]
The fourth line, after a newline with ten points, and some more words so that it runs on.

Case B, in a quote:
\begin{quote}\small\ttfamily
The first line of the quote.\\[4pt]
The second line, after 4pt.\\
The third line.\\[4pt]
The fourth line, after 4pt.
\end{quote}
Case C, space in the middle of a paragraph: a sentence long enough to fill most of a line or two of the text.
\vspace{6pt} And after the vspace another sentence, which the space goes below the line of. Then a few more words
follow so that the paragraph has lines to break at every width. A lemma states something.\smallskip{} Its next
sentence comes after a small skip, as in real papers, and the paragraph goes on for another line or so.

Case D, space above a line: a paragraph with words enough to fill two lines of the text at every width it is set at,
and here\vadjust pre{\vskip 5pt} the line this word is on gets five points above it. It goes on a little further.

Case E, material that draws: words enough for a line or two at every width, and then a box\vadjust{\kern2pt\hbox{\itshape set below the line this word ends up in}}
goes below this line, and here a rule\vadjust{\kern3pt\hrule\kern3pt} across the whole column below its line. And a box\vadjust pre{\hbox{\itshape set above the line}} goes above this line, and the paragraph goes on.

The end, a plain paragraph after all of them.
\end{document}
